% $Id$ %
\subsection{OTP Client}
This plugin provides the ability to generate one-time-passwords for
authentication purposes. It implements an HMAC-based One-Time Password
Algorithm (RFC 4226), and on targets which support it, a Time-Based
One-Time Password Algorithm (RFC 6238).

\opt{rtc}{ It is important to note that for TOTP accounts to work
  properly, the clock on your device MUST be accurate to no less than
  30 seconds from the time on the authentication server. See
  \reference{ref:Timeanddateactual} for more information.  }

\subsubsection{Adding Accounts}
The plugin supports two methods of adding accounts: URI import, and
manual entry.

\subsubsection{URI Import}
This method of adding an account reads a list of URIs from a file. It
expects each URI to be on a line by itself in the following format:

\begin{verbatim}
otpauth://[hotp OR totp]/[account name]?secret=[Base32 secret][&counter=X][&period=X][&digits=X]
\end{verbatim}

An example is shown below, provisioning a TOTP key for an account called ``bob'':

\begin{verbatim}
otpauth://totp/bob?secret=JBSWY3DPEHPK3PXP
\end{verbatim}

Any other options are not supported and will be ignored.

Most services will provide a scannable QR code that encodes a OTP
URI. In order to use those, first scan the QR code and save the URI to
a file on your device. If necessary, rewrite the URI so it is in the
format shown above.

\subsection{Manual Import}
If URI import is not possible, the plugin supports the manual entry of
data associated with an account. While some services provide you with
a base32-encoded string to use as the shared secret, others may
not. In this case, you can scan the QR code provided and enter secret
parameter as provided. Be careful not to make a mistake while entering
the Base32 secret for an account.